\chapter{Definición del MVP}

Este capítulo define formalmente el Producto Mínimo Viable (MVP) de ContractIA:
qué construimos, para quién, qué problema resuelve, qué queda fuera del alcance
y cómo sabemos que el MVP cumplió su propósito. La definición está alineada con
el autodiagnóstico del Capítulo~5 y sirve de referencia para la validación
presentada en el Capítulo~9.

% ─────────────────────────────────────────────────────────────────────────────
\section{Propuesta de Valor}

ContractIA entrega una propuesta de valor concreta y medible:

\begin{quote}
\textit{``Un especialista legal de ProInversión puede obtener, en aproximadamente
1 hora y sin asistencia del equipo técnico, un reporte completo de inconsistencias
contractuales con trazabilidad exacta a cláusula, evidencia textual y referencia
normativa, reduciendo el esfuerzo de screening inicial de 320 horas a menos del
1\% del tiempo original.''}
\end{quote}

Esta propuesta se diferencia de las soluciones existentes en tres dimensiones:

\begin{table}[H]
\centering
\caption{Diferenciación de ContractIA frente a enfoques alternativos}
\label{tab:diferenciacion}
\begin{tabular}{p{3.5cm} p{3.5cm} p{3.5cm} p{3.5cm}}
\toprule
\textbf{Dimensión} & \textbf{Revisión manual} &
\textbf{Herramientas de búsqueda} & \textbf{ContractIA MVP} \\
\midrule
Tiempo de screening & 320 horas & 200 horas & $\approx$1 hora \\[3pt]
Cobertura & $\sim$60\% (muestreo) & $\sim$70\% & 100\% \\[3pt]
Detección semántica & Alta (humano) & Baja & Alta (LLM) \\[3pt]
Conocimiento normativo & Requiere experto & No integrado & RAG integrado \\[3pt]
Trazabilidad & Manual / parcial & Parcial & Automática 100\% \\[3pt]
Costo estimado & \$7{,}200/contrato & \$4{,}000/contrato & \$4--8/contrato \\
\bottomrule
\end{tabular}
\end{table}

% ─────────────────────────────────────────────────────────────────────────────
\section{User Stories Priorizadas (MoSCoW)}

Las funcionalidades del MVP fueron priorizadas aplicando el marco MoSCoW
(Must Have / Should Have / Could Have / Won't Have). La Tabla~\ref{tab:moscow}
resume las historias de usuario que guiaron el desarrollo.

\begin{table}[H]
\centering
\caption{Priorización MoSCoW de User Stories del MVP}
\label{tab:moscow}
\begin{tabular}{p{0.8cm} p{4.5cm} p{3.5cm} p{4.5cm}}
\toprule
\textbf{Prior.} & \textbf{Como\ldots} & \textbf{Quiero\ldots} &
\textbf{Para\ldots} \\
\midrule
\textbf{M} & Especialista legal & Cargar un contrato PDF y recibir un
reporte de inconsistencias en $\approx$1 hora & Reducir el tiempo de
screening inicial de 320 horas \\[4pt]
\textbf{M} & Especialista legal & Ver la ubicación exacta (capítulo,
sección, cláusula) de cada hallazgo & Verificar y priorizar la revisión
especializada \\[4pt]
\textbf{M} & Especialista legal & Que cada hallazgo incluya la referencia
normativa que lo sustenta & Fundamentar observaciones sin búsqueda manual
de normas \\[4pt]
\textbf{M} & Especialista legal & Dejar el análisis corriendo y recibirlo
por email al terminar & No tener que esperar frente a la pantalla durante
el procesamiento \\[4pt]
\textbf{S} & Auditor senior & Ver un resumen ejecutivo con hallazgos
clasificados por severidad & Priorizar la atención en riesgos ALTA antes
de revisar los MEDIA/BAJA \\[4pt]
\textbf{S} & Administrador & Gestionar usuarios y aprobar accesos desde
el bot Telegram & Controlar quién puede usar el sistema sin exponer la API \\[4pt]
\textbf{C} & Especialista legal & Comparar dos versiones del mismo contrato
y ver qué cambió & Identificar modificaciones entre adendas sin revisión
manual \\[4pt]
\textbf{W} & Institución & Conectar ContractIA con el repositorio
documental oficial & (Fase 2 del roadmap; requiere integración
institucional) \\
\bottomrule
\end{tabular}
\end{table}

\noindent\textbf{Leyenda:} M = Must Have (imprescindible en MVP),
S = Should Have (importante, incluido si el tiempo lo permite),
C = Could Have (deseable, postergado), W = Won't Have (explícitamente
fuera del alcance de este MVP).

% ─────────────────────────────────────────────────────────────────────────────
\section{Alcance Funcional}

El MVP implementa las siguientes capacidades funcionales:

\begin{enumerate}
    \item \textbf{Ingesta y OCR:} Procesamiento de contratos PDF (y DOCX,
    aunque no recomendado) con extracción de texto, normalización UTF-8
    y detección de páginas ilegibles. Sin límite de páginas establecido;
    el tiempo de procesamiento escala proporcionalmente a la extensión
    del documento.
    \item \textbf{Segmentación automática:} Detección de capítulos y anexos
    mediante expresiones regulares con mecanismo anti-TOC y resolución de
    duplicados por densidad de cláusulas.
    \item \textbf{Indexación multinivel:} Construcción de tres índices
    (Secciones, Global de Cláusulas, Local por sección).
    \item \textbf{GraphRAG:} Construcción de grafo de conocimiento contractual
    con extracción de tripletas (entidad, relación, entidad) por sección.
    \item \textbf{RAG híbrido:} Recuperación de contexto normativo mediante
    FAISS + BM25 con reranking por Cohere.
    \item \textbf{Auditoría multi-agente:} Ejecución paralela de tres agentes
    (Jurista, Auditor, Cronista) sobre cada sección.
    \item \textbf{Clasificación de hallazgos:} Severidad automática (ALTA /
    MEDIA / BAJA) con trazabilidad a cláusula y párrafo.
    \item \textbf{Reportes:} Generación de informe ejecutivo y reporte técnico
    en PDF, con grafo de conocimiento visualizado.
    \item \textbf{API REST:} Endpoints autenticados (JWT) para integración con
    sistemas externos.
    \item \textbf{Interfaces de usuario:} Frontend web (Next.js) y bot
    Telegram con flujo de carga, cola de trabajos (máx.\ 5 por usuario),
    procesamiento en segundo plano y notificación por email al finalizar.
\end{enumerate}

\section{Alcance No Funcional}

\begin{itemize}
    \item Tiempo de procesamiento: aproximadamente 1 hora para contratos
    de concesión de gran extensión (300+ páginas) con Gemini~2.5~Pro;
    escala proporcionalmente según volumen de secciones.
    \item Disponibilidad: $\geq 99\%$ en Cloud Run (auto-scaling).
    \item Autenticación: JWT con roles (usuario, auditor, admin).
    \item Privacidad: Cumplimiento con Ley N$^\circ$~29733 y estándares GCP.
\end{itemize}

\section{Criterios de Éxito}

Los criterios de éxito del MVP están alineados con los objetivos SMART
definidos en el Capítulo~3. La Tabla~\ref{tab:criterios_exito} resume
los indicadores, valores objetivo y resultados obtenidos en la validación.

\begin{table}[H]
\centering
\caption{Criterios de éxito del MVP ContractIA}
\label{tab:criterios_exito}
\begin{tabular}{p{4cm} p{3cm} p{3cm} p{3.5cm}}
\toprule
\textbf{Criterio} & \textbf{Métrica} & \textbf{Objetivo} &
\textbf{Resultado (validación)} \\
\midrule
Tiempo de screening & TTR (horas) & $\approx 1$ hora &
  $\approx 1$h $15$min $\checkmark$ \\[4pt]
Cobertura de análisis & \% cláusulas auditadas & $> 95\%$ &
  $100\%$ $\checkmark$ \\[4pt]
Precisión semántica & \% hallazgos válidos & $\geq 85\%$ &
  $87\%$ $\checkmark$ \\[4pt]
Recall & \% inconsistencias detectadas & $\geq 85\%$ &
  $94\%$ $\checkmark$ \\[4pt]
Trazabilidad & \% hallazgos con ubicación exacta & $100\%$ &
  $100\%$ $\checkmark$ \\[4pt]
Flujo autónomo & Sistema operable sin equipo técnico & Sí &
  Sí $\checkmark$ \\
\bottomrule
\end{tabular}
\end{table}

% ─────────────────────────────────────────────────────────────────────────────
\section{Exclusiones}

Las siguientes capacidades fueron deliberadamente excluidas del MVP para
mantener el foco en la hipótesis central y entregar valor demostrable
en el tiempo del proyecto:

\begin{itemize}
    \item \textbf{Procesamiento simultáneo entre usuarios:} El sistema
    implementa una cola individual por usuario (máx.\ 5 trabajos pendientes)
    con un contrato activo a la vez por usuario; no hay paralelismo entre
    usuarios simultáneos ni procesamiento en GPU.
    \item \textbf{Integración con sistemas documentales institucionales:}
    No hay conexión con SharePoint, SEACE, INFOBRAS ni repositorios
    gubernamentales; la carga de contratos es manual vía interfaz web o bot.
    \item \textbf{Módulo predictivo de riesgo contractual longitudinal:}
    No se incluye análisis histórico entre contratos ni clasificación
    predictiva de riesgo; está planificado en la Fase 4 del roadmap.
    \item \textbf{Contratos en formato DOCX:} Soportados técnicamente pero
    no recomendados en producción; el sistema está optimizado para PDF.
    \item \textbf{Autenticación SSO institucional:} El acceso se gestiona
    mediante lista blanca de usuarios en el bot Telegram y JWT en la API;
    no hay integración con Active Directory ni proveedores de identidad
    corporativos.
    \item \textbf{Comparación de versiones de contratos:} La detección
    de diferencias entre adendas o versiones del mismo contrato está
    planificada en la Fase 2 del roadmap.
\end{itemize}

% ─────────────────────────────────────────────────────────────────────────────
\section{Criterios de Aceptación del MVP}

El MVP se considera entregado y funcional cuando cumple simultáneamente
las siguientes condiciones verificables:

\begin{enumerate}
    \item Un usuario con acceso aprobado puede cargar un contrato PDF
    desde el bot Telegram o la interfaz web sin intervención del equipo
    de desarrollo.
    \item El sistema procesa el contrato de forma autónoma (segmentación,
    indexación, RAG, análisis multi-agente) y genera el reporte sin
    errores de pipeline.
    \item El usuario recibe notificación por email al finalizar el
    análisis y puede descargar el reporte desde la interfaz.
    \item El reporte contiene hallazgos clasificados por severidad
    (ALTA / MEDIA / BAJA) con ubicación exacta (capítulo, sección,
    cláusula, párrafo) y referencia normativa cuando aplica.
    \item El tiempo total de procesamiento para un contrato de referencia
    (324 páginas, 41 secciones) no supera las 2 horas en condiciones
    normales de operación.
    \item El sistema permanece disponible y operable sin reinicio manual
    durante al menos 7 días consecutivos de uso.
\end{enumerate}